\chapter{Duckietown System Integration (\texttt{patos})}

\section{Objective}

This chapter covers the engineering extension that brings the planning stack
into an executable Duckietown workflow. The goal is robust closed-loop operation,
not only offline trajectory generation.

\section{Runtime architecture}

The runtime is organized into three coordinated modules:

\begin{enumerate}
    \item manager node (state aggregation and acceleration publication),
    \item Julia TCP bridge (receding-horizon solver service),
    \item per-robot adapter node (acceleration-to-command conversion).
\end{enumerate}

\begin{figure}[tbp]
    \centering
    \fbox{\parbox{0.92\linewidth}{
    \textbf{Placeholder figure: \texttt{patos} runtime architecture.}
    Include ROS topics and TCP request/response path between manager and Julia
    bridge.
    }}
    \caption{Runtime architecture of the \texttt{patos} integration.}
    \label{fig:patos_runtime_arch}
\end{figure}

\section{Manager behavior}

The manager supports multi-robot operation and can run in:

\begin{itemize}
    \item trajectory-sequence mode (CSV-driven),
    \item Julia-solver mode (online solver queries).
\end{itemize}

In solver mode, it sends per-robot $(x,y,v_x,v_y,\theta)$ snapshots and
receives $(a_x,a_y)$ commands. Error handling is explicit: on solver timeout or
failure, the manager can publish zero acceleration or fallback controls depending
on configuration.

\section{Adapter behavior}

Each adapter converts incoming acceleration vectors into bounded
velocity/yaw-rate commands using an internal velocity-vector controller with
state decay and heading correction. The adapter also publishes telemetry,
including model state, command state, and optional encoder/wheel-command
feedback.

This makes the controller chain observable and debuggable for both simulation
and physical runs.

\section{Julia bridge details}

The bridge maintains persistent context for receding-horizon progression and
provides simple commands (e.g., \texttt{PING}, \texttt{STEP}, \texttt{RESET}).
Internally it builds and solves a constrained multi-player planning problem with
road, speed, and collision constraints plus goal-seeking objectives.

\section{Deployment workflow}

The repository includes scripts/UI for:

\begin{itemize}
    \item container build and launch,
    \item code transfer to one or multiple robots,
    \item starting/stopping manager and adapters,
    \item telemetry inspection and map-level visualization.
\end{itemize}

\section{Validation path}

Validation is layered:

\begin{enumerate}
    \item unit-level controller tests,
    \item local simulation checks,
    \item ROS container smoke tests,
    \item real robot runs.
\end{enumerate}

This sequence reduces integration risk and supports repeatable debugging.

\section{Current maturity}

The system is functional and can run end-to-end. Remaining work for final thesis
results is focused on broader scenario coverage and quantitative robustness
metrics.

\section{Positioning}

This work extends simulation-first GOOP literature by providing a concrete
systems bridge to an embodied platform \cite{duckietown:2017,delasheras:2025}.
